\section*{Chapter \ref{ch:g12:inflation-and-money} --- Inflation and the Quantity of Money}
% figdata: figdata/grade-12/03-inflation-and-money/out/ (qtm, qtm_key, velocity); Kestria's numbers invented.

\begin{solution}{exo:g12:inflation-and-money:1}
$P = MV/Y = 200 \times 5/500 = 2$.
\end{solution}

\begin{solution}{exo:g12:inflation-and-money:2}
$V = 3000/600 = 5$. Prices double: $PY = MV$ doubles with $Y$ fixed.
\end{solution}

\begin{solution}{exo:g12:inflation-and-money:3}
About $8 - 3 = 5\,\%$.
\end{solution}

\begin{solution}{exo:g12:inflation-and-money:4}
Japan lies below the line: inflation (0.5\,\%) well under money growth (2.0\,\%), so output grew or velocity fell, or both. The Democratic Republic of the Congo lies above
it: inflation (173\,\%) exceeded money growth (149\,\%), so velocity rose or output shrank; people held less money as it lost value.
\end{solution}

\begin{solution}{exo:g12:inflation-and-money:5}
Below: 116 of the 123 countries have lower inflation than money growth. By the growth-rate form, inflation is money growth plus velocity growth minus output growth; output
grows in most countries, and velocity tends to fall as people and firms hold more of their \omterm{def:g11:inequality:wealth}{wealth} in bank deposits.
\end{solution}

\begin{solution}{exo:g12:inflation-and-money:6}
2007: $14\,474/7294 \approx 1.98$. 2019: $21\,373/14\,837 \approx 1.44$.
\end{solution}

\begin{solution}{exo:g12:inflation-and-money:7}
$V = 2400/800 = 3$. Prices rise in proportion to money: by 25\,\%.
\end{solution}

\begin{solution}{exo:g12:inflation-and-money:8}
Exactly $1.30/1.02 - 1 \approx 27.5\,\%$; approximately $30 - 2 = 28\,\%$.
\end{solution}

\begin{solution}{exo:g12:inflation-and-money:9}
$0.05 \times 40 = 2$ a month.
\end{solution}

\begin{solution}{exo:g12:inflation-and-money:10}
Exactly $1.07/1.04 - 1 \approx 2.88\,\%$; approximately $7 - 4 = 3\,\%$.
\end{solution}

\begin{solution}{exo:g12:inflation-and-money:11}
About 1.78 in 1960; highest about 2.19, in 1997; 1.18 in 2020.
\end{solution}

\begin{solution}{exo:g12:inflation-and-money:12}
In logarithms over the twelve years, taking the consumer price index for $P$: money $\ln 2.03 \approx 0.71$, prices $\ln 1.24 \approx 0.21$, velocity
$\ln(1.44/1.98) \approx -0.32$. The identity then gives real output growth of $0.71 - 0.32 - 0.21 \approx 0.18$, about 19\,\%. Nothing contradicts the identity, but the theory's assumption of stable velocity failed: half the money growth was
absorbed by a fall in velocity, as people and banks held more money at low interest rates.
\end{solution}

\begin{solution}{exo:g12:inflation-and-money:13}
$0.2/1.2 \approx 16.7\,\%$; $1/2 = 50\,\%$; $10/11 \approx 90.9\,\%$. The rate tends to 1: inflation can take at most all the money people hold. Since people hold less and
less money as inflation rises, the revenue, the rate times real balances, eventually falls however fast money is printed.
\end{solution}

\begin{solution}{exo:g12:inflation-and-money:14}
The repayment is 1050. Expected real value $1050/1.02 \approx 1029.4$; actual $1050/1.10 \approx 954.5$, a real rate of about $-4.5\,\%$ instead of about $+2.9\,\%$. The
borrower gains about 75 in real terms, at the lender's expense.
\end{solution}

\begin{solution}{exo:g12:inflation-and-money:15}
The money was printed to pay for deficits: in Germany in May 1922 taxes covered only 21\,\% of the government's income, and the central bank bought the rest of its debt. A
new currency printed for the same deficit would have gone the same way. The German reform paired the Rentenmark with an end to the government's borrowing at the central bank
and planned cuts in its staff; Zimbabwe's ended when it stopped printing its own currency at all.
\end{solution}

\begin{solution}{pb:g12:inflation-and-money:1}
\textbf{1.} $m = Y/V = 100/2 = 50$.

\textbf{2.} $\mu \times 50 = 10$, so $\mu = 20\,\%$ a month.

\textbf{3.} $(1.2)(1)/(1) - 1 = 20\,\%$ a month.

\textbf{4.} $0.2/1.2 \approx 16.7\,\%$.

\textbf{5.} No: 20\,\% is below 50\,\% a month.

\textbf{6.} $(1.2)(2)/1 - 1 = 140\,\%$ in that month.

\textbf{7.} $m = 100/4 = 25$.

\textbf{8.} $0.2 \times 25 = 5$ a month: half as much.

\textbf{9.} $\mu \times 25 = 10$ gives $\mu = 40\,\%$, and inflation of 40\,\% a month at constant velocity.

\textbf{10.} Faster money growth raises expected inflation; people hold still less money, velocity rises again, real balances shrink, and the same revenue needs faster
money growth still.

\textbf{11.} None: $\mu = 0$.

\textbf{12.} $(1)(1)/(1) - 1 = 0$.

\textbf{13.} Velocity halves with money and output fixed, so prices halve: a fall of 50\,\%.

\textbf{14.} Real balances can rise from 25 to 50 at unchanged prices: 25 of new money, in real terms.

\textbf{15.} People hold more money only if they expect it to keep its value; a new currency issued under new rules marks the break and lowers expected inflation at once,
so velocity falls back sooner.

\textbf{16.} $30 \times 0.2/1.2 = 5$ in real terms.

\textbf{17.} Its real value is divided by 2.4: from 1200 to 500 in prices of the start of the month, a fall of about 58\,\%.

\textbf{18.} It too falls by about 58\,\%: the worker pays for the inflation that the firm gains from.

\textbf{19.} Money growth of 20\,\% a month dwarfs the changes in output and, outside the month of flight, in velocity; at low inflation, velocity moves as much as money.

\textbf{20.} In the month velocity doubles, inflation is 140\,\%. At money growth of 20\,\% a month, \omterm{def:g12:inflation-and-money:seigniorage}{seigniorage} is 10 a month before and 5 a month after.
\end{solution}
